\section{Background}
\label{sec:background}
% ============================================================

A coding agent combines an LLM with tools for reading files, editing code, executing shell commands, and interacting with external services. Representative systems include Devin AI~\cite{devin2024}, Claude Code~\cite{claudecode2025}, and OpenHands~\cite{wang2025openhands}. During each step, the model receives an observation, proposes an action, and receives the executor's result~\cite{yao2022react}. Observations may include user instructions, repository content, tool output, and error traces. Actions may include tool calls, shell commands, and code edits. The loop ends when the task is complete or when it reaches a limit on steps, cost, or time. Because one session can contain many tool calls, a security policy must remain active throughout the trajectory.

We model a coding agent as $A = (M, \mathcal{T}, P)$, where $M$ is the LLM, $\mathcal{T} = \{t_1, \ldots, t_k\}$ is the tool set, and $P$ is the system prompt. At step $i$, the agent observes history $h_{<i}$, generates action $a_i = M(P \oplus h_{<i})$, and receives observation $o_i = \mathcal{T}(a_i)$ from the executor. The operator controls $P$ and includes it in every model call. Under our threat model, requests cannot modify $P$, although users may infer its contents. The system prompt can therefore carry a policy throughout the session, but its finite size limits how much policy text can be included.

\emph{Skills} are modular natural-language documents that extend or specialize agent behavior without retraining~\cite{anthropic2025skills}. The emerging Agent Skills standard~\cite{agentskills2025} packages each skill as a \texttt{SKILL.md} file with metadata and an instruction body. Its progressive-disclosure protocol initially loads a skill's name and description, then lets the model load the body when a request appears relevant. We use the document format because it makes policies easy to inspect and update. We do not use progressive disclosure for security policy. A malicious first request could influence the model before it loads the policy or could induce the model to skip loading it. \sysname therefore places the complete security skill in $P$ before the session begins. This design gives the policy coverage from the first step at the cost of prompt space on every request.
